% ============================================================================
%  Capitolo: Sistema di Monitoraggio Real-Time per Applicazioni DDS
%  Autore: Leonardo Leggeri
%  Data: Aprile 2026
% ============================================================================

\chapter{Sistema di Monitoraggio Real-Time per Applicazioni DDS}
\label{ch:monitoraggio}

Questo capitolo descrive la progettazione, l'implementazione e il funzionamento
del sistema di monitoraggio real-time sviluppato per analizzare il comportamento
temporale di applicazioni concorrenti che utilizzano il middleware
DDS (Data Distribution Service) per la comunicazione inter-processo.

Il sistema e' composto da tre componenti principali:
\begin{enumerate}
    \item Un'\textbf{applicazione C++ strumentata} (\texttt{GeometryBroadcastner})
          con due thread periodici applicativi, che emette marker di tracciamento
          nel kernel Linux tramite ftrace. Il middleware eProsima FastDDS crea
          automaticamente i propri thread interni per la comunicazione.
    \item Un \textbf{analizzatore C++} (\texttt{thread\_analysis}) che elabora
          i dati grezzi del trace e produce un file CSV strutturato.
    \item Un \textbf{visualizzatore Python} (\texttt{monitorRealTime.py}) che
          genera diagrammi temporali (Gantt chart) con statistiche dettagliate.
\end{enumerate}

L'intera catena e' orchestrata dallo script \texttt{run\_trace.sh}, che automatizza
l'acquisizione, l'analisi e la visualizzazione in un unico comando.


% ────────────────────────────────────────────────────────────────────────────
\section{Architettura dell'Applicazione Strumentata}
\label{sec:architettura}

L'applicazione \texttt{GeometryBroadcastner} e' strutturata con
\textbf{due thread periodici applicativi} e sfrutta il thread interno del
middleware eProsima FastDDS per la comunicazione sulla rete.
La struttura dei thread applicativi e' derivata dal pattern dell'applicazione
di riferimento \texttt{app10h}.

% ............................................................................
\subsection{I Thread dell'Applicazione}
\label{subsec:thread}

\begin{table}[htbp]
\centering
\caption{Thread visibili nel diagramma temporale.}
\label{tab:thread}
\begin{tabular}{llccp{5.5cm}}
\toprule
\textbf{Thread} & \textbf{Tipo} & \textbf{Periodo} & \textbf{Priorita'} & \textbf{Funzione} \\
\midrule
\texttt{DDS\_Publish} & Periodico & 600\,ms & 39 (RM) &
    Esegue il carico computazionale, prepara il messaggio \texttt{Point}
    e chiama \texttt{publish()} sul middleware DDS. \\
\texttt{Activity\_1}  & Periodico & 250\,ms & 74 (RM) &
    Carico computazionale puro, identico ad \texttt{Activity\_1}
    nell'applicazione di riferimento \texttt{app10h}. \\
\midrule
\texttt{DDS\_Comm}    & Event-driven & --- & 120 (default) &
    Thread \textbf{reale} del middleware eProsima FastDDS
    (\texttt{dds.ev.0}). Viene creato automaticamente dalla libreria
    FastDDS al momento dell'inizializzazione del \texttt{DomainParticipant}.
    Gestisce l'invio effettivo dei messaggi sulla rete. \\
\bottomrule
\end{tabular}
\end{table}

Il punto chiave dell'architettura e' che \texttt{DDS\_Comm} \textbf{non e'
un thread creato dall'applicazione}: e' il vero thread evento
(\texttt{dds.ev.0}) del middleware eProsima FastDDS. Quando il task
periodico \texttt{DDS\_Publish} chiama \texttt{publish()}, il middleware
serializza i dati e sveglia il proprio thread evento interno per gestire
l'invio sulla rete. L'analizzatore \texttt{thread\_analysis} lo individua
automaticamente nel trace tramite il nome OS e lo rinomina
\texttt{DDS\_Comm} per chiarezza nella visualizzazione.

Questa architettura permette di:
\begin{itemize}
    \item \textbf{Osservare} il comportamento reale del middleware DDS
          senza modificarlo o strumentarlo.
    \item \textbf{Visualizzare} sulla timeline esattamente quando il
          thread interno di FastDDS si attiva per inviare il messaggio.
    \item \textbf{Correlare} l'istante in cui l'applicazione chiama
          \texttt{publish()} (marker \texttt{DDS\_MSG\_START} sulla corsia
          \texttt{DDS\_Publish}) con l'istante in cui il middleware
          effettivamente lavora (blocchi RUN sulla corsia \texttt{DDS\_Comm}).
    \item \textbf{Analizzare} eventuali ritardi tra la richiesta di invio
          e l'effettiva trasmissione.
\end{itemize}

% ............................................................................
\subsection{Politica di Scheduling}
\label{subsec:scheduling}

I due thread applicativi utilizzano la politica \texttt{SCHED\_FIFO}
(scheduling real-time POSIX) e sono vincolati alla \textbf{CPU~0} tramite
\texttt{pthread\_attr\_setaffinity\_np}. Questo vincolo e' intenzionale:
forzando entrambi i thread sullo stesso core, si creano le condizioni per
osservare preemption e competizione per il tempo macchina.

Le priorita' sono assegnate secondo il principio del
\textbf{Rate Monotonic Scheduling} (RMS):
\begin{equation}
    \text{priorita'} = 99 - \left\lfloor \frac{\text{periodo\_ms}}{10} \right\rfloor
    \label{eq:rm_priority}
\end{equation}

Applicando l'Equazione~\ref{eq:rm_priority}:
\begin{itemize}
    \item \texttt{Activity\_1} (periodo 250\,ms): priorita' $99 - 25 = 74$
    \item \texttt{DDS\_Publish} (periodo 600\,ms): priorita' $99 - 60 = 39$
\end{itemize}

Il thread \texttt{DDS\_Comm} (\texttt{dds.ev.0}) e' creato internamente
dal middleware FastDDS con priorita' di default del sistema operativo
(SCHED\_OTHER, nice~0). Non avendo vincoli real-time, il suo scheduling
dipende interamente dal kernel Linux. Questo lo rende particolarmente
interessante da osservare: eventuali ritardi nell'invio dei messaggi
possono essere causati dalla preemption da parte dei task real-time
applicativi.

% ............................................................................
\subsection{Flusso della Comunicazione DDS}
\label{subsec:dds_flow}

La comunicazione tra l'applicazione e la rete avviene attraverso il
middleware eProsima FastDDS. Il flusso e' il seguente:

\begin{enumerate}
    \item Il thread \texttt{DDS\_Publish} esegue il carico computazionale
          (\texttt{activity\_load}).
    \item Prepara il messaggio \texttt{Point} con le coordinate aggiornate.
    \item Scrive il marker \texttt{DDS\_MSG\_START\_DDS\_Publish} su ftrace.
    \item Chiama \texttt{g\_broadcastner.publish(\&point\_msg)}.
          Internamente, FastDDS:
          \begin{enumerate}
              \item Serializza il messaggio nel formato CDR.
              \item Inserisce i dati nella history cache del \texttt{DataWriter}.
              \item Sveglia il proprio thread evento (\texttt{dds.ev.0})
                    che gestisce l'invio effettivo sulla rete.
          \end{enumerate}
    \item La chiamata \texttt{publish()} ritorna al thread applicativo.
    \item Scrive il marker \texttt{DDS\_MSG\_END\_DDS\_Publish}.
\end{enumerate}

\begin{lstlisting}[language=C++, caption={Funzione del task periodico DDS\_Publish.}, label={lst:publish_func}]
void dds_publish_function(void *instance) {
    /* carico computazionale */
    activity_load(10);

    /* prepara il messaggio */
    Point point_msg;
    point_msg.x(1.0 * g_counter);
    point_msg.y(2.0 * g_counter);
    point_msg.z(3.0 * g_counter);

    /* comunicazione DDS con marker */
    write_trace_marker("DDS_MSG_START_DDS_Publish");
    g_broadcastner.publish(&point_msg);
    write_trace_marker("DDS_MSG_END_DDS_Publish");
}
\end{lstlisting}

I marker \texttt{DDS\_MSG\_START} e \texttt{DDS\_MSG\_END} delimitano
la chiamata \texttt{publish()} dal punto di vista del thread applicativo.
Nel frattempo, il thread interno \texttt{dds.ev.0} (visualizzato come
\texttt{DDS\_Comm}) si attiva in modo asincrono per completare l'invio.
Confrontando i timestamp dei marker con i blocchi RUN di \texttt{DDS\_Comm}
e' possibile misurare il ritardo tra la richiesta applicativa e l'effettiva
trasmissione.


% ────────────────────────────────────────────────────────────────────────────
\section{Strumentazione con Ftrace}
\label{sec:ftrace}

Il sistema di tracciamento si basa su \textbf{ftrace}, il framework di
tracing integrato nel kernel Linux. Ftrace permette di registrare eventi
di scheduling (\texttt{sched\_switch}, \texttt{sched\_wakeup}) e marker
applicativi scritti direttamente dal codice utente.

% ............................................................................
\subsection{Marker Applicativi}
\label{subsec:marker}

I marker vengono scritti tramite il file speciale
\texttt{/sys/kernel/tracing/trace\_marker}. Ogni scrittura su questo file
viene registrata da ftrace con il timestamp, il PID del thread che ha
scritto e il testo del marker.

\begin{lstlisting}[language=C, caption={Funzione per scrivere marker ftrace.}, label={lst:write_marker}]
static inline void write_trace_marker(const char *msg) {
    static int trace_fd = -1;
    if (trace_fd < 0) {
        trace_fd = open("/sys/kernel/tracing/trace_marker", O_WRONLY);
    }
    if (trace_fd >= 0) {
        write(trace_fd, msg, strlen(msg));
    }
}
\end{lstlisting}

L'applicazione emette i seguenti marker, suddivisi per thread:

\begin{table}[htbp]
\centering
\caption{Marker ftrace emessi dall'applicazione.}
\label{tab:marker}
\begin{tabular}{llp{6cm}}
\toprule
\textbf{Thread emittente} & \textbf{Marker} & \textbf{Significato} \\
\midrule
\texttt{DDS\_Publish} & \texttt{PERIOD\_START\_DDS\_Publish} & Inizio del periodo \\
                       & \texttt{FUNCTION\_START\_DDS\_Publish} & Inizio esecuzione funzione \\
                       & \texttt{DDS\_MSG\_START\_DDS\_Publish} & Inizio chiamata \texttt{publish()} \\
                       & \texttt{DDS\_MSG\_END\_DDS\_Publish} & Ritorno dalla \texttt{publish()} \\
                       & \texttt{FUNCTION\_END\_DDS\_Publish} & Fine esecuzione funzione \\
                       & \texttt{SLEEP\_START\_DDS\_Publish} & Inizio attesa prossimo periodo \\
                       & \texttt{PERIOD\_END\_DDS\_Publish} & Fine del periodo \\
\midrule
\texttt{Activity\_1}  & \texttt{PERIOD\_START\_Activity\_1} & Inizio del periodo \\
                       & \texttt{FUNCTION\_START\_Activity\_1} & Inizio esecuzione funzione \\
                       & \texttt{FUNCTION\_END\_Activity\_1} & Fine esecuzione funzione \\
                       & \texttt{SLEEP\_START\_Activity\_1} & Inizio attesa prossimo periodo \\
                       & \texttt{PERIOD\_END\_Activity\_1} & Fine del periodo \\
\bottomrule
\end{tabular}
\end{table}

\paragraph{Punto chiave.}
I marker \texttt{DDS\_MSG\_START} e \texttt{DDS\_MSG\_END} sono scritti
dal thread \texttt{DDS\_Publish} e delimitano la chiamata
\texttt{publish()} dal punto di vista dell'applicazione.
Il thread reale del middleware (\texttt{dds.ev.0}, visualizzato come
\texttt{DDS\_Comm}) \textbf{non scrive marker}: la sua attivita' e'
visibile unicamente tramite gli eventi di scheduling del kernel
(\texttt{sched\_switch}), che l'analizzatore registra come blocchi
RUN, SLEEP e PREEMPT sulla sua corsia nel diagramma.

Nel trace, i marker appaiono cosi':
\begin{verbatim}
DDS_Publish-8270 [000] 1132.50012: print: tracing_mark_write: DDS_MSG_START_DDS_Publish
DDS_Publish-8270 [000] 1132.50145: print: tracing_mark_write: DDS_MSG_END_DDS_Publish
\end{verbatim}

E il thread reale FastDDS appare negli eventi di scheduling:
\begin{verbatim}
    <idle>-0     [007] 1132.50015: sched_switch: swapper/7:0 R ==> dds.ev.0:11614
  dds.ev.0-11614 [007] 1132.50018: sched_switch: dds.ev.0:11614 S ==> swapper/7:0
\end{verbatim}

% ............................................................................
\subsection{Naming dei Thread con \texttt{pthread\_setname\_np}}
\label{subsec:naming}

I thread applicativi vengono nominati tramite la funzione GNU
\texttt{pthread\_setname\_np}, sia dall'interno del thread stesso
(tramite la libreria \texttt{activity\_library}) sia dal main dopo
la creazione:

\begin{lstlisting}[language=C, caption={Assegnazione del nome ai thread applicativi.}, label={lst:setname}]
// Nella libreria PeriodicTask:
pthread_setname_np(pthread_self(), activity.name);

// Nel main, dopo pthread_create:
pthread_setname_np(thread1, "DDS_Publish");
pthread_setname_np(thread2, "Activity_1");
\end{lstlisting}

Il thread \texttt{dds.ev.0} e' nominato automaticamente dal middleware
eProsima FastDDS al momento della creazione del \texttt{DomainParticipant}.
Il nome \texttt{dds.ev.0} segue la convenzione di FastDDS dove:
\begin{itemize}
    \item \texttt{dds.} e' il prefisso comune a tutti i thread del middleware.
    \item \texttt{ev} indica il thread evento (event loop).
    \item \texttt{0} e' l'indice del \texttt{DomainParticipant}.
\end{itemize}

Nel trace, i nomi OS appaiono nel campo \texttt{comm}:
\begin{verbatim}
     DDS_Publish-8270  [000]  ...    (thread applicativo, nome impostato da noi)
     Activity_1-8271   [000]  ...    (thread applicativo, nome impostato da noi)
     dds.ev.0-11614    [007]  ...    (thread eProsima FastDDS, nome automatico)
\end{verbatim}


% ────────────────────────────────────────────────────────────────────────────
\section{Il Task Periodico: Struttura e Funzionamento}
\label{sec:periodic_task}

I thread \texttt{DDS\_Publish} e \texttt{Activity\_1} eseguono la stessa
funzione generica \texttt{PeriodicTask}, definita nella libreria condivisa
\texttt{activity\_library.c}. Questa funzione implementa un loop periodico
con le seguenti caratteristiche:

\begin{itemize}
    \item \textbf{Attivazione periodica} tramite \texttt{clock\_nanosleep}
          con \texttt{CLOCK\_MONOTONIC} e \texttt{TIMER\_ABSTIME}.
    \item \textbf{Rilevamento deadline miss}: se il tempo di fine esecuzione
          supera il tempo di rilascio del prossimo periodo, il job corrente
          e' marcato come miss e il job successivo viene saltato.
    \item \textbf{Marker ftrace completi}: ogni fase del periodo
          (inizio, esecuzione, fine, sleep) e' delimitata da marker.
\end{itemize}

\begin{lstlisting}[language=C, caption={Struttura dati per i parametri di attivita'.}, label={lst:activity_par}]
typedef struct activity_parameters {
    char name[50];
    void (*function)(void* instance);
    int period;
    long int deadline;
    void* instance;
    bool print;
} t_activity_par;
\end{lstlisting}

Il loop periodico segue questo flusso per ogni periodo:

\begin{enumerate}
    \item Scrive \texttt{PERIOD\_START\_<nome>}.
    \item Calcola il tempo di rilascio del prossimo periodo:
          \texttt{time\_add\_millisecs(\&release, period)}.
    \item Esegue la funzione applicativa (delimitata da
          \texttt{FUNCTION\_START} e \texttt{FUNCTION\_END}).
    \item Controlla se \texttt{exec\_end\_time > next\_release\_time}
          (deadline miss).
    \item Scrive \texttt{SLEEP\_START\_<nome>}.
    \item Dorme fino al prossimo rilascio con \texttt{clock\_nanosleep}.
    \item Scrive \texttt{PERIOD\_END\_<nome>}.
\end{enumerate}


% ────────────────────────────────────────────────────────────────────────────
\section{Pipeline di Acquisizione e Analisi}
\label{sec:pipeline}

L'intera pipeline di acquisizione, analisi e visualizzazione e' automatizzata
dallo script \texttt{run\_trace.sh}. Il flusso e' composto da cinque fasi.

\begin{figure}[htbp]
\centering
\begin{tikzpicture}[
    node distance=1.8cm,
    box/.style={draw, rounded corners, minimum width=3.5cm,
                minimum height=0.9cm, align=center, font=\small},
    arrow/.style={->, thick, >=stealth}
]
\node[box] (app) {Applicazione C++\\(\texttt{GeometryBroadcastner})};
\node[box, right=of app] (trace) {\texttt{trace-cmd record}\\(kernel ftrace)};
\node[box, right=of trace] (report) {\texttt{trace-cmd report}\\(\texttt{trace\_output.txt})};
\node[box, below=1.2cm of report] (analysis) {\texttt{thread\_analysis}\\(\texttt{timeline.csv})};
\node[box, left=of analysis] (python) {\texttt{monitorRealTime.py}\\(grafico PNG)};

\draw[arrow] (app) -- (trace) node[midway, above, font=\scriptsize] {ftrace marker};
\draw[arrow] (trace) -- (report) node[midway, above, font=\scriptsize] {\texttt{trace.dat}};
\draw[arrow] (report) -- (analysis) node[midway, right, font=\scriptsize] {testo};
\draw[arrow] (analysis) -- (python) node[midway, above, font=\scriptsize] {CSV};
\end{tikzpicture}
\caption{Pipeline di acquisizione e analisi del trace.}
\label{fig:pipeline}
\end{figure}

% ............................................................................
\subsection{Fase 1: Acquisizione con \texttt{trace-cmd}}
\label{subsec:tracecmd}

Lo script esegue l'applicazione sotto \texttt{trace-cmd record}, che
configura ftrace per catturare gli eventi \texttt{sched:sched\_switch}
(cambio di contesto) e \texttt{sched:sched\_wakeup} (risveglio di un
thread), oltre ai marker applicativi scritti su \texttt{trace\_marker}:

\begin{lstlisting}[language=bash, caption={Comando di acquisizione.}, label={lst:tracecmd}]
sudo trace-cmd record \
    -e sched:sched_switch \
    -e sched:sched_wakeup \
    -o "$OUTPUT_DIR/trace.dat" \
    "$EXECUTABLE"
\end{lstlisting}

Il risultato e' un file binario \texttt{trace.dat} che contiene tutti
gli eventi registrati durante l'esecuzione.

% ............................................................................
\subsection{Fase 2: Conversione in Testo}
\label{subsec:report}

Il file binario viene convertito in formato testuale leggibile:

\begin{lstlisting}[language=bash]
trace-cmd report trace.dat > trace_output.txt
\end{lstlisting}

Ogni riga del file testuale ha il formato:
\begin{verbatim}
<comm>-<PID> [<CPU>] <timestamp>: <evento>: <dettagli>
\end{verbatim}

Ad esempio:
\begin{verbatim}
DDS_Publish-8270 [000] 1131.997662: print: tracing_mark_write: PERIOD_START_DDS_Publish
Activity_1-8271  [000] 1131.998123: sched_switch: Activity_1:8271 [74] S ==> DDS_Publish:8270 [39]
DDS_Publish-8270 [000] 1132.500123: print: tracing_mark_write: DDS_MSG_START_DDS_Publish
\end{verbatim}


% ────────────────────────────────────────────────────────────────────────────
\section{L'Analizzatore: \texttt{thread\_analysis}}
\label{sec:thread_analysis}

L'analizzatore e' un programma C++ che legge il file di trace testuale
e il file sorgente dell'applicazione, ed esporta un file
\texttt{timeline.csv} contenente intervalli di stato, marker e statistiche.

% ............................................................................
\subsection{Fase 1: Estrazione dei Parametri dal Sorgente}
\label{subsec:source_parsing}

L'analizzatore legge il file sorgente C/C++ dell'applicazione per estrarre
automaticamente i nomi dei task, i periodi e le deadline. Utilizza tre
espressioni regolari:

\begin{lstlisting}[language=C++, caption={Regex per estrazione parametri dal sorgente.}, label={lst:regex}]
// Nome del task
regex name_re(R"(sprintf\s*\(\s*(\w+)\.name\s*,\s*"([^"]+)"\s*\))");
// Periodo in millisecondi
regex period_re(R"((\w+)\.period\s*=\s*(\d+))");
// Deadline in millisecondi
regex deadline_re(R"((\w+)\.deadline\s*=\s*(\d+))");
\end{lstlisting}

Applicando queste regex al sorgente di \texttt{GeometryBroadcastner.cpp},
l'analizzatore estrae:

\begin{table}[htbp]
\centering
\caption{Parametri estratti automaticamente dal sorgente.}
\label{tab:params}
\begin{tabular}{lccc}
\toprule
\textbf{Variabile} & \textbf{Nome} & \textbf{Periodo (ms)} & \textbf{Deadline (ms)} \\
\midrule
\texttt{activity\_1} & \texttt{DDS\_Publish} & 600 & 600 \\
\texttt{activity\_2} & \texttt{Activity\_1}  & 250 & 250 \\
\bottomrule
\end{tabular}
\end{table}

Il thread \texttt{DDS\_Comm} (\texttt{dds.ev.0}) non appare in questa
tabella perche' non e' definito nel sorgente dell'applicazione: e' un
thread interno del middleware eProsima FastDDS. Viene scoperto nella fase
successiva tramite il suo nome OS nel trace.

% ............................................................................
\subsection{Fase 2: Discovery Dinamica dei Thread}
\label{subsec:discovery}

L'analizzatore effettua due passate sul file di trace per scoprire
tutti i thread da tracciare.

\paragraph{Passata 1: Scoperta dei thread applicativi via marker.}
Scorre tutte le righe del trace cercando eventi di tipo \texttt{print}
o \texttt{tracing\_mark\_write}. Per ogni PID che scrive un marker,
registra la coppia (PID, nome\_OS):

\begin{lstlisting}[language=C++, caption={Scoperta dei thread applicativi tramite marker.}, label={lst:discovery1}]
if (event.find("print") != std::string::npos ||
    event.find("tracing_mark_write") != std::string::npos) {
    if (marker_pids.count(pid) == 0)
        marker_pids[pid] = comm;  // es: marker_pids[8270] = "DDS_Publish"
}
\end{lstlisting}

\paragraph{Scoperta del thread evento FastDDS.}
Nella stessa passata, l'analizzatore cerca specificamente i thread il
cui nome OS inizia con \texttt{dds.ev} --- il thread evento di eProsima
FastDDS responsabile dell'invio dei messaggi. Quando lo trova, lo
registra con il nome \texttt{DDS\_Comm}:

\begin{lstlisting}[language=C++, caption={Scoperta e rinomina del thread reale FastDDS.}, label={lst:discovery_ev}]
// Tracciamo il thread evento di eProsima (dds.ev.*) che gestisce
// l'invio dei messaggi. Lo rinominiamo "DDS_Comm" nel CSV.
if (tasks.count(pid) == 0 && comm.find("dds.ev") == 0) {
    tasks[pid].name = "DDS_Comm";
    tasks[pid].expected_period_sec = 0;   // non periodico
    tasks[pid].expected_deadline_sec = 0;
}
\end{lstlisting}

\paragraph{Passata 2: Registrazione con filtro.}
I PID che hanno scritto marker vengono registrati come task, con un
filtro che esclude tutti gli altri thread interni del middleware FastDDS
(\texttt{dds.shm.*}, \texttt{dds.udp.*}, \texttt{dds.asyn.*}, etc.):

\begin{lstlisting}[language=C++, caption={Filtro dei thread interni FastDDS non rilevanti.}, label={lst:filter}]
// Ignora thread interni FastDDS che scrivono marker
// (es. dds.dsha scrive DDS_READ_START/END per il listener)
if (comm.find("dds.") == 0 || comm.find("tpool") == 0 ||
    comm.find("fastrtps") == 0 || comm.find("fastdds") == 0) {
    continue;  // non tracciare
}
\end{lstlisting}

Questo filtro e' fondamentale: il middleware FastDDS crea internamente
numerosi thread (\texttt{dds.shm.*} per la shared memory,
\texttt{dds.udp.*} per il trasporto UDP, \texttt{dds.asyn.*} per il
flow control, etc.) che inquinerebbero il diagramma temporale.
Solo il thread evento \texttt{dds.ev.0}, responsabile dell'effettivo
invio dei messaggi, viene tracciato e rinominato \texttt{DDS\_Comm}.

% ............................................................................
\subsection{Fase 3: Costruzione degli Intervalli}
\label{subsec:intervals}

Nella terza passata, l'analizzatore raccoglie gli eventi \texttt{sched\_switch}
e i marker per ogni PID registrato, e costruisce tre tipi di intervalli:

\begin{description}
    \item[RUN] Il thread e' in esecuzione sulla CPU.
        Costruito accoppiando gli eventi \texttt{switch\_in}
        (il thread viene schedulato) con \texttt{switch\_out}
        (il thread cede la CPU).
    \item[SLEEP] Il thread e' in attesa volontaria (stato \texttt{S} o \texttt{D}).
        Corrisponde ai periodi in cui il thread dorme
        (\texttt{clock\_nanosleep}, \texttt{pthread\_cond\_wait}, etc.).
    \item[PREEMPT] Il thread e' pronto ma non in esecuzione (stato \texttt{R}).
        Un thread a priorita' superiore lo ha interrotto.
\end{description}

% ............................................................................
\subsection{Fase 4: Costruzione degli Intervalli DDS\_MSG}
\label{subsec:dds_msg_intervals}

Dopo aver estratto tutti i marker, l'analizzatore cerca le coppie
\texttt{DDS\_MSG\_START} / \texttt{DDS\_MSG\_END} per ogni task e le
accoppia in ordine cronologico:

\begin{lstlisting}[language=C++, caption={Accoppiamento dei marker DDS\_MSG.}, label={lst:dds_pairing}]
// Per ogni DDS_MSG_START, trova il primo DDS_MSG_END successivo
while (di < dds_starts.size() && dj < dds_ends.size()) {
    if (dds_ends[dj] > dds_starts[di]) {
        double ds = dds_starts[di];
        double de = dds_ends[dj];
        double dur_ms = (de - ds) * 1000.0;
        csv << tdata.name << ",DDS_MSG," << ds << "," << de
            << "," << dur_ms << "\n";
        di++; dj++;
    } else {
        dj++;
    }
}
\end{lstlisting}

Ogni coppia genera una riga nel CSV di tipo \texttt{DDS\_MSG} con
timestamp di inizio, fine e durata in millisecondi. Vengono inoltre
calcolate tre statistiche:

\begin{itemize}
    \item \texttt{STAT\_DDS\_COUNT}: numero totale di comunicazioni DDS.
    \item \texttt{STAT\_DDS\_WCET}: Worst-Case Execution Time della
          comunicazione (durata massima).
    \item \texttt{STAT\_DDS\_ACET}: Average-Case Execution Time
          (durata media).
\end{itemize}

% ............................................................................
\subsection{Fase 5: Statistiche per Task Periodici}
\label{subsec:stats}

Per i task periodici, l'analizzatore calcola un insieme completo di
metriche real-time:

\begin{table}[htbp]
\centering
\caption{Statistiche calcolate per ogni task periodico.}
\label{tab:stats}
\begin{tabular}{lp{8cm}}
\toprule
\textbf{Metrica} & \textbf{Descrizione} \\
\midrule
Periodo misurato & Mediana delle distanze tra marker \texttt{PERIOD\_START}
                   consecutivi. \\
Jitter & Deviazione standard delle distanze tra \texttt{PERIOD\_START}. \\
WCET & Worst-Case Execution Time: massimo tempo di esecuzione (RUN)
       osservato in un singolo periodo. \\
ACET & Average-Case Execution Time: media del tempo di esecuzione. \\
WCRT & Worst-Case Response Time: massimo tempo tra \texttt{PERIOD\_START}
       e \texttt{FUNCTION\_END}. \\
Utilizzazione & Rapporto tra tempo totale di RUN e tempo totale
                osservato. \\
Deadline miss & Numero di periodi in cui il \texttt{FUNCTION\_END}
               supera la deadline. \\
Worst slack & Peggiore margine tra deadline e tempo di risposta
              (negativo = miss). \\
\bottomrule
\end{tabular}
\end{table}


% ────────────────────────────────────────────────────────────────────────────
\section{Il Formato CSV}
\label{sec:csv}

Il file \texttt{timeline.csv} e' il contratto tra l'analizzatore C++ e
il visualizzatore Python. Contiene tre categorie di righe:

\paragraph{Righe di stato} (intervalli temporali):
\begin{verbatim}
task,type,start,end,duration_ms
DDS_Publish,RUN,1131.997662,1132.290601,292.939
DDS_Publish,SLEEP,1132.290601,1132.797944,507.343
DDS_Publish,DDS_MSG,1132.500123,1132.501456,1.333
DDS_Comm,RUN,1132.500150,1132.500180,0.030
Activity_1,PREEMPT,1132.013012,1132.013138,0.126
\end{verbatim}

\paragraph{Righe di marker} (eventi puntuali):
\begin{verbatim}
DDS_Publish,MARKER_PERIOD_START_DDS_Publish,1131.997662,1131.997662,0.0
\end{verbatim}

\paragraph{Righe di statistica}:
\begin{verbatim}
DDS_Publish,STAT_PERIOD,0,0,600.000000
DDS_Publish,STAT_WCET,0,0,292.939000
DDS_Publish,STAT_DDS_COUNT,0,0,6
DDS_Publish,STAT_DDS_WCET,0,0,1.500000
\end{verbatim}


% ────────────────────────────────────────────────────────────────────────────
\section{Il Visualizzatore: \texttt{monitorRealTime.py}}
\label{sec:visualizer}

Il visualizzatore Python legge il \texttt{timeline.csv} e genera un
diagramma temporale (Gantt chart) con le seguenti caratteristiche.

% ............................................................................
\subsection{Layout del Grafico}
\label{subsec:layout}

Il grafico e' suddiviso verticalmente in piu' pannelli:

\begin{enumerate}
    \item \textbf{Timeline} (pannello superiore): diagramma temporale con
          una corsia orizzontale per ogni thread. Ogni corsia mostra i
          blocchi di stato (RUN, PREEMPT, SLEEP) e gli overlay DDS\_MSG.
    \item \textbf{Deadline Miss Log} (pannello centrale, se presente):
          tabella dettagliata delle deadline miss.
    \item \textbf{Task Statistics} (pannello inferiore): tabella con tutte
          le metriche calcolate per ogni task.
\end{enumerate}

% ............................................................................
\subsection{Rendering degli Intervalli DDS\_MSG}
\label{subsec:dds_rendering}

Gli intervalli \texttt{DDS\_MSG} vengono visualizzati come
\textbf{blocchi arancioni} (\texttt{\#FF6F00}) sovrapposti ai blocchi
di stato nella corsia del thread \texttt{DDS\_Publish} (dove la
\texttt{publish()} viene chiamata). I blocchi sono:

\begin{itemize}
    \item Leggermente piu' sottili dei blocchi RUN (38\% dell'altezza)
          e spostati verso l'alto (offset 28\%), in modo da essere
          visivamente distinguibili.
    \item Delimitati da \textbf{linee verticali} (pin) agli estremi,
          per evidenziare con precisione l'istante di inizio e fine
          della comunicazione.
    \item Annotati con la \textbf{durata in millisecondi} sopra il blocco,
          se sufficientemente larghi.
\end{itemize}

% ............................................................................
\subsection{Codifica Cromatica e Legenda}
\label{subsec:colori}

\begin{table}[htbp]
\centering
\caption{Codifica cromatica degli elementi del diagramma temporale.}
\label{tab:colori}
\begin{tabular}{llp{6cm}}
\toprule
\textbf{Colore} & \textbf{Elemento} & \textbf{Significato} \\
\midrule
\textcolor[HTML]{43A047}{Verde}  & RUN & Il thread e' in esecuzione sulla CPU. \\
\textcolor[HTML]{E53935}{Rosso}  & PREEMPT & Il thread e' pronto ma e' stato
    interrotto da un thread a priorita' superiore. \\
\textcolor[HTML]{90A4AE}{Grigio} & SLEEP & Il thread e' in attesa volontaria. \\
\textcolor[HTML]{FF6F00}{Arancione} & DDS\_MSG & Intervallo della chiamata
    \texttt{publish()} sul thread applicativo (dal momento in cui
    l'applicazione invoca la funzione al suo ritorno). \\
Nero (solido)  & Fine Periodo & Linea verticale che segna la fine
    di ogni periodo del task. \\
Rosso (tratteggiato) & Deadline & Linea verticale alla scadenza della
    deadline (se diversa dal periodo). \\
$\blacktriangledown$ Rosso & Deadline Miss & Triangolo che indica un
    superamento della deadline. \\
\bottomrule
\end{tabular}
\end{table}


% ────────────────────────────────────────────────────────────────────────────
\section{Il Build System}
\label{sec:build}

Il sistema utilizza un \texttt{Makefile} unificato che gestisce la
compilazione di tutte le applicazioni (locali e DDS) e automatizza
la copia dei file necessari al tool di tracing.

% ............................................................................
\subsection{Struttura delle Cartelle}
\label{subsec:folders}

\begin{verbatim}
parser/
  |-- Makefile               # Build system unificato
  |-- run_trace.sh           # Script di orchestrazione
  |-- thread_analysis.cpp    # Analizzatore C++
  |-- monitorRealTime.py     # Visualizzatore Python
  |-- sorgenti/              # File sorgente (copiati dal post-build)
  |-- eseguibili/            # Eseguibili (copiati dal post-build)
  |-- output/traces/         # Risultati dei trace
  |-- dds/
       |-- lib/              # Librerie DDS (time, activity, communication)
       |-- msg/              # Definizioni IDL dei messaggi
       |-- app/
            |-- GeometryPeriodic/   # Applicazione principale
\end{verbatim}

% ............................................................................
\subsection{Regola Post-Build}
\label{subsec:postbuild}

Dopo la compilazione, il Makefile copia automaticamente tutti gli
eseguibili nella cartella \texttt{eseguibili/} e tutti i file sorgente
nella cartella \texttt{sorgenti/}. Questa organizzazione e' necessaria
affinche' lo script \texttt{run\_trace.sh} possa:

\begin{enumerate}
    \item Trovare l'eseguibile da tracciare in un percorso standard.
    \item Trovare il file sorgente corrispondente per passarlo
          all'analizzatore, che ne estrae i parametri di periodo e deadline.
\end{enumerate}

Il comando per compilare tutto e':
\begin{lstlisting}[language=bash]
make -j4
\end{lstlisting}


% ────────────────────────────────────────────────────────────────────────────
\section{Utilizzo del Sistema}
\label{sec:utilizzo}

Per eseguire un trace completo dell'applicazione \texttt{GeometryBroadcastner}:

\begin{lstlisting}[language=bash]
# 1. Compilazione
make -j4

# 2. Acquisizione trace (richiede privilegi root per ftrace)
sudo ./run_trace.sh GeometryBroadcastner
\end{lstlisting}

Lo script esegue automaticamente tutte le fasi della pipeline e produce:
\begin{itemize}
    \item \texttt{output/traces/trace\_results\_GeometryBroadcastner\_<timestamp>/}
    \begin{itemize}
        \item \texttt{trace.dat} --- dati grezzi del trace
        \item \texttt{trace\_output.txt} --- report testuale
        \item \texttt{timeline.csv} --- dati strutturati
        \item \texttt{risultati\_finali.txt} --- statistiche a terminale
        \item \texttt{timeline\_monitor.png} --- diagramma temporale
    \end{itemize}
\end{itemize}

Le cartelle di output vengono create con i permessi dell'utente corrente
(tramite \texttt{chown} post-acquisizione), in modo da poter essere
eliminate senza privilegi di root.


% ────────────────────────────────────────────────────────────────────────────
\section{Interpretazione del Diagramma Temporale}
\label{sec:interpretazione}

Il diagramma temporale prodotto dal sistema permette di rispondere
alle seguenti domande analitiche:

\begin{enumerate}
    \item \textbf{Quando l'applicazione richiede l'invio?}
          I blocchi arancioni \texttt{DDS\_MSG} sulla corsia
          \texttt{DDS\_Publish} mostrano l'intervallo della chiamata
          \texttt{publish()} dal punto di vista del thread applicativo.
    \item \textbf{Quando il middleware effettivamente invia?}
          I blocchi verdi RUN sulla corsia \texttt{DDS\_Comm}
          (\texttt{dds.ev.0}) mostrano quando il thread reale di
          eProsima FastDDS si attiva per trasmettere il messaggio.
    \item \textbf{Qual e' il ritardo tra richiesta e invio?}
          Confrontando il timestamp di \texttt{DDS\_MSG\_START}
          (corsia \texttt{DDS\_Publish}) con il primo blocco RUN
          successivo sulla corsia \texttt{DDS\_Comm}, si misura la
          latenza introdotta dal middleware.
    \item \textbf{C'e' interferenza tra comunicazione e computazione?}
          Blocchi rossi (PREEMPT) sulla corsia \texttt{DDS\_Comm}
          indicano che il thread del middleware e' stato interrotto
          dai task real-time applicativi a priorita' superiore.
    \item \textbf{Il task periodico rispetta le deadline?}
          Le linee verticali nere (fine periodo) e i triangoli rossi
          (deadline miss) forniscono una risposta immediata.
    \item \textbf{Qual e' il jitter del task periodico?}
          La tabella delle statistiche riporta il jitter misurato
          come deviazione standard delle distanze tra attivazioni.
\end{enumerate}
